\scanpage{268}
% Faint reversed arrow in the upper blank area shows through from the other side.

\mspage{- 83 -}\hfill IX

7.1.3. La projectivité de $\hat B$ signifie, en vertu de EGA III 5.4.5, que l'on peut trouver un Module inversible $\underline{\hat M}$ sur $\hat B$, i.e. une suite de Modules inversibles $\underline M_n$ sur les $B_n$ qui se "recollen\deleted{\unsure{t"}}\unsure{t"} pour $n$ variable, qui soit \uline{ample} relativement à $S$ i.e. qui soit tel que $\underline M_{\unsure{o|0}}$ soit ample. L'idée naturelle serait de partir d'un Module inversible ample $\underline M$ de $G$, et de montrer que pour tout $n$, la restriction $\underline L_n$ de $\underline L$ à $G_n$ provient par image inverse d'un Module inversible $M_n$ sur $B_n$. J'ignore si de tels $M_n$ existent nécessairement (et en doute fort). Nous allons suivre une méthode un peu différente, en construisant une \uline{polarisation} de $B$, i.e. un système projectif de \uline{polarisations} $\rho_n$ des $B_n$ qui se recollent. On sait alors [5] que les doubles de ces polarisations définissent canoniquement des Modules inversibles $\underline M_n$ (symétriques) sur les $B_n$, qui par leur nature canonique se recolleront nécessairement, et fourniront le $\underline M$ cherché.

Nous allons interpréter un élément de $\mathrm{NS}(B_n/S_n)$ comme une correspondance divisorielle symétrique sur $(B_n,B_n)$ [5], i.e. comme une biextension de $(B_n,B_n)$ par $\mathbb G_{m\,S_n}$. En vertu de VIII 3.5, ces biextensions correspondent biunivoquement aux biextensions de $(A_n,A_n)$ par $\mathbb G_{mS}$. Pour obtenir un système cohérent de telles biextensions $W_n$ il suffit de partir d'une

\[
W,\ \text{biextension symétrique de }(A,A)\text{ par }\mathbb G_{mS}\,,\tag*{(7.1.4)}
\]

en prenant les $W_n$ se déduisant de $W$ par les changements de base $S_n\longrightarrow S$. Nous devons de plus supposer que l'élément de $\mathrm{NS}(B_o)$ défini par $W_o$ est une polarisation de $B_o$ pour obtenir ainsi une polarisation de $\hat B$.

% archive sheet number: 267
% stamp: 401 (lower centre, on label); faint partly covered earlier number above.
